\chapter{First-Order Partial Derivatives and the Gradient}

When a function depends on more than one input, we measure how it responds
to a change in \emph{one} variable at a time, holding the others fixed,
the ceteris paribus principle. This is exactly what a partial derivative
computes. From these one-variable slopes we assemble the gradient, which
packages all first-order information into a single vector and unlocks the
geometry of multivariate optimisation.

\section{Core Rules}

\begin{rulebox}{Topological vocabulary of domains}
\begin{itemize}[leftmargin=1.6em,itemsep=0.3ex]
  \item The \emph{open ball} (or \emph{neighbourhood}) of radius $r>0$ around $\mathbf{x}_0\in\R^n$ is
        $B_r(\mathbf{x}_0)=\{\mathbf{x}\in\R^n :
        \|\mathbf{x}-\mathbf{x}_0\|<r\}$.
  \item $\mathbf{x}_0\in M$ is an \emph{interior point} of $M$ if some open ball around
        $\mathbf{x}_0$ lies entirely in $M$.
  \item $M$ is \emph{open} if \emph{every} point of $M$ is
        an interior point.
  \item $\mathbf{x}_0$ is a \emph{boundary point} of
        $M$ if every ball around $\mathbf{x}_0$ contains points of $M$
        \emph{and} points of its complement.
  \item $M$ is \emph{closed} if it contains all
        its boundary points; equivalently, the complement
        $\R^n\setminus M$ is open.
  \item $M$ is \emph{bounded} if
        $M\subseteq B_r(\mathbf{0})$ for some radius $r$.
  \item $M\subseteq\R^n$ is \emph{compact} exactly when
        $M$ is closed \emph{and} bounded.
\end{itemize}
Economic examples: the open positive orthant $\{(x,y): x>0,\ y>0\}$
(the natural Cobb--Douglas domain) is open and unbounded; a budget set
$\{\mathbf{x}\in\R^n : \mathbf{p}\cdot\mathbf{x}\le m,\
\mathbf{x}\ge\mathbf{0}\}$ is closed and bounded, hence compact.
\end{rulebox}

\begin{rulebox}{Why it matters}
Differentiability (and hence every partial derivative below) is
defined at \emph{interior} points of the domain, where $f$ can be
probed in all directions. And by the \emph{Weierstrass theorem}, a
continuous function on a compact set attains its maximum and minimum:
this is why budget-constrained optimisation problems are
guaranteed to have solutions.
\end{rulebox}

\begin{rulebox}{Definition}
Let $f:\R^n\to\R$. The partial derivative of $f$ with respect to $x_i$ at
$\mathbf{x}=(x_1,\ldots,x_n)$ is
\[
  \pd{f}{x_i}(\mathbf{x}) \;=\; \lim_{h\to0}
    \frac{f(x_1,\ldots,x_i+h,\ldots,x_n)-f(x_1,\ldots,x_n)}{h} .
\]
\end{rulebox}

\begin{rulebox}{Practical computation}
To compute $\partial f/\partial x_i$, differentiate $f$ with respect to
$x_i$ and treat every other variable as a constant. All single-variable
rules (power, product, quotient, chain) apply unchanged.
\end{rulebox}

\begin{rulebox}{Notation}
The following all refer to the same object:
\[
  \pd{f}{x} = f_x = \partial_x f = D_x f .
\]
\end{rulebox}

\begin{rulebox}{Gradient}
The vector of all first-order partial derivatives is
\[
  \grad f(\mathbf{x}) \;=\;
  \begin{pmatrix} f_{x_1}(\mathbf{x}) \\ f_{x_2}(\mathbf{x}) \\ \vdots \\
    f_{x_n}(\mathbf{x}) \end{pmatrix} .
\]
\end{rulebox}

\begin{rulebox}{Geometric meaning}
\begin{itemize}[leftmargin=1.6em,itemsep=0.3ex]
  \item $\grad f(\mathbf{x}_0)$ points in the direction of \emph{steepest
        ascent} at $\mathbf{x}_0$.
  \item $\|\grad f(\mathbf{x}_0)\|$ is the maximal rate of increase at that
        point.
  \item $\grad f(\mathbf{x}_0)$ is \emph{perpendicular} to the level set
        $\{\mathbf{x} : f(\mathbf{x})=f(\mathbf{x}_0)\}$ at $\mathbf{x}_0$.
\end{itemize}
\end{rulebox}

\begin{rulebox}{Economic interpretation}
\begin{itemize}[leftmargin=1.6em,itemsep=0.3ex]
  \item For a production function $Q=F(K,L)$: $F_K=\mathrm{MP}_K$
        (marginal product of capital), $F_L=\mathrm{MP}_L$ (marginal
        product of labour).
  \item For a utility function $U(x,y)$: $U_x, U_y$ are the marginal
        utilities; the marginal rate of substitution is
        $\mathrm{MRS}_{xy}=U_x/U_y$.
  \item For production: the marginal rate of technical substitution is
        $\mathrm{MRTS}_{LK}=F_L/F_K$.
\end{itemize}
\end{rulebox}

\begin{rulebox}{Directional derivative}
Let $\mathbf{u}\in\R^n$ be a \emph{unit} vector ($\|\mathbf{u}\|=1$). The
directional derivative of $f$ at $\mathbf{x}_0$ in the direction
$\mathbf{u}$ is
\[
  D_{\mathbf{u}}f(\mathbf{x}_0) \;=\; \grad f(\mathbf{x}_0)\cdot\mathbf{u} .
\]
\end{rulebox}

\begin{rulebox}{Caution}
If you are given a non-unit direction $\mathbf{v}$, you \emph{must}
normalise first: $\mathbf{u}=\mathbf{v}/\|\mathbf{v}\|$. Forgetting this is
the single most common error in directional-derivative problems.
\end{rulebox}

\section{Practice Problems}

\begin{problems}
  \item Compute $f_x$ and $f_y$ for $f(x,y)=5x^4-2x^2y^3+y^5$.
  \item For $F(K,L)=12\,K^{1/4}L^{3/4}$, find $F_K$ and $F_L$ at
        $(K,L)=(16,81)$.
  \item Compute $f_x$ and $f_y$ for $f(x,y)=x^2e^{3y}$.
  \item Compute $\nabla f$ for $f(x,y)=4x^3y-xy^4+7$.
  \item Compute the directional derivative of $f(x,y)=x^2-4xy+y^3$ at
        $(2,1)$ in the direction $\mathbf{v}=(5,-12)$.
  \item For $F(K,L)=6K^{2/3}L^{1/3}$, compute $\mathrm{MP}_K$ and
        $\mathrm{MP}_L$ at $(K,L)=(27,8)$.
  \item Compute $f_x$ and $f_y$ for $f(x,y)=(3x-2y)^4$.
  \item Compute $f_x$ and $f_y$ for $f(x,y)=e^{xy^2}$.
  \item For $F(K,L)=25K^{0.6}L^{0.5}$, find $F_K$ and $F_L$.
  \item Compute $\nabla f$ for $f(x,y,z)=x^2+3yz-z^2e^x$.
  \item Compute the directional derivative of $f(x,y)=e^x\ln(1+y)$ at
        $(0,0)$ in the direction $\mathbf{v}=(1,2)$.
  \item For $U(x,y)=2\ln(x)+5\ln(y)$, derive $\mathrm{MRS}_{xy}=U_x/U_y$ as
        a function of $(x,y)$ and evaluate at $(3,5)$.
  \item Simplify and compute $f_x, f_y$ for
        $f(x,y)=\dfrac{x^2-y^2}{x+y}$.
  \item Compute $f_x$ and $f_y$ for $f(x,y)=\ln(x^2y+1)$.
  \item For $U(x_1,x_2,x_3)=x_1^{1/2}x_2^{1/3}x_3^{1/6}$, compute the
        three partials.
  \item Compute the directional derivative of $f(x,y,z)=xy-2yz+z^2$ at
        $(1,2,3)$ in the direction $\mathbf{v}=(1,2,2)$.
  \item For $F(K,L)=K^{0.25}L^{0.75}$, derive $\mathrm{MRTS}_{LK}=F_L/F_K$
        and evaluate at $(K,L)=(16,4)$.
  \item Expand and then compute $f_x, f_y$ for
        $f(x,y)=(x^2+y)(x-3y^2)$.
  \item Compute $f_x$ and $f_y$ for $f(x,y)=xe^{-y/x}$ with $x\neq0$.
  \item For $F(K,L)=(3K^{1/2}+2L^{1/2})^2$, find $F_K$ and $F_L$.
  \item Compute $\nabla f$ for $f(x,y)=\arctan(y/x)$ with $x>0$.
  \item For $f(x,y)=x^{1/2}y^{1/2}$ at $(K,L)=(4,9)$, find the unit
        direction of steepest ascent and the maximum rate of increase.
  \item For the CES function $F(K,L)=(K^{1/2}+L^{1/2})^2$, verify Euler's
        theorem $KF_K+LF_L=rF$ for the appropriate degree of homogeneity
        $r$.
  \item Compute $f_x$ and $f_y$ for $f(x,y)=\dfrac{2x+3y}{x^2+y^2}$.
  \item Compute $f_x$ and $f_y$ for $f(x,y)=\dfrac{e^x}{1+e^y}$.
  \item For $F(K,L)=A(\alpha K^\rho+(1-\alpha)L^\rho)^{\nu/\rho}$ (CES),
        find $F_L$.
  \item Compute the directional derivative of $f(x,y)=\ln(x^2+y^2)$ at
        $(3,4)$ in the direction $\mathbf{v}=(-1,1)$.
  \item Compute $f_x, f_y, f_z$ for $f(x,y,z)=x^2yz^3-4xy^2z+5z$.
  \item Compute $f_x$ and $f_y$ for $f(x,y)=\sqrt{4x^2+9y^2}$.
  \item For $F(K,L)=(K^{1/2}+L^{1/2})^4$, compute $F_K$ at
        $(K,L)=(4,9)$.
  \item Compute $\nabla f$ for $f(x,y,z)=\ln(x^2+y^2+z^2)$.
  \item Compute $f_x$ and $f_y$ for
        $f(x,y)=\ln\left(\sqrt{x^2+y^2}\right)$.
  \item Cost minimisation. A firm uses inputs $(x,y)$ with cost
        $C(x,y)=4x+9y$ and output $Q(x,y)=x^{1/2}y^{1/2}$. Compute
        $\nabla C$ and $\nabla Q$, and verify that at the cost-minimising
        input mix $\nabla C=\lambda\nabla Q$ for some $\lambda$.
  \item Compute $f_x$ and $f_y$ for $f(x,y)=(1+xy)^5$.
  \item Compute the directional derivative of $f(x,y)=x^3-2xy+y^2$ at
        $(1,2)$ in the direction $\mathbf{v}=(3,4)$.
  \item For $F(K,L)=10K^{0.3}L^{0.7}$ (Cobb--Douglas), compute $F_K$ and
        $F_L$ at $(K,L)=(100,100)$ and verify $KF_K+LF_L=F$.
  \item Compute $f_x$ and $f_y$ for $f(x,y)=\sin(xy)+\cos(x-y)$.
  \item Compute $\nabla f$ for $f(x,y)=x^2y-xy^2+3x-2y$ at the point
        $(1,-1)$.
  \item For $U(x,y)=x^\alpha y^{1-\alpha}$ (Cobb--Douglas,
        $0<\alpha<1$), derive $\mathrm{MRS}_{xy}=U_x/U_y$ and show it
        equals $\dfrac{\alpha}{1-\alpha}\cdot\dfrac{y}{x}$.
  \item Compute the directional derivative of $f(x,y)=e^{xy}$ at $(1,0)$
        in the direction $\mathbf{v}=(-3,4)$.
  \item Compute $f_x$ and $f_y$ for $f(x,y)=x^y$ (with $x>0$).
        \emph{Hint:} write $x^y=e^{y\ln x}$.
  \item For $F(K,L)=\min\{2K,3L\}$ (Leontief, fixed-proportions, on the
        smooth region away from the kink), compute $F_K$ and $F_L$ in
        each of the two regions $\{2K<3L\}$ and $\{2K>3L\}$.
  \item Compute $\nabla f$ for $f(x,y,z)=xyz+x^2-z^2$ at the point
        $(2,-1,3)$.
  \item Compute the unit vector of steepest \emph{descent} for
        $f(x,y)=x^2+4y^2$ at $(3,1)$, and the corresponding rate of
        decrease.
  \item For $f(x,y)=xe^y-ye^x$, compute $f_x, f_y$.
  \item Compute $f_x$ and $f_y$ for $f(x,y)=\dfrac{x}{x^2+y^2}$.
  \item For $U(x,y)=x+2\sqrt{y}$ (quasi-linear), derive
        $\mathrm{MRS}_{xy}$.
  \item Compute the directional derivative of $f(x,y,z)=x^2+y^2-z^2$ at
        $(2,-1,2)$ in the direction $\mathbf{v}=(2,2,1)$.
  \item Compute $f_x$ and $f_y$ for $f(x,y)=\tan(x+y)+\tan(x-y)$.
  \item For $F(K,L)=(\alpha K+\beta L)^r$ (linear-aggregate,
        $\alpha,\beta,r>0$), compute $F_K, F_L$, and verify
        $KF_K+LF_L=rF$.
  \item Compute $\nabla f$ for $f(x,y)=e^{-x}\sin(y)+e^x\cos(y)$.
  \item Compute the directional derivative of $f(x,y)=\sqrt{x^2+y^2}$ at
        $(3,4)$ in the direction $\mathbf{v}=(1,0)$.
  \item For $U(x,y)=\min\{2x,3y\}$, compute the marginal utilities on
        each smooth region $\{2x<3y\}$ and $\{2x>3y\}$.
  \item Compute $f_x, f_y$ for $f(x,y)=\dfrac{\ln(x+y)}{x-y}$ (assume
        $x>y>0$).
  \item For $F(K,L)=AK^\alpha L^\beta$ general Cobb--Douglas, derive
        $\mathrm{MRTS}_{LK}$ and show it depends only on the ratio $K/L$
        (not on scale).
  \item Compute the unit direction of steepest ascent for
        $f(x,y)=x^2-y^2$ at $(2,-1)$, and the maximum rate.
  \item Compute $f_x$ and $f_y$ for $f(x,y)=\cosh(x+y)+\sinh(x-y)$.
  \item For $f(x,y,z)=e^{x+y+z}$, compute $\nabla f$ at $(0,0,0)$ and the
        directional derivative in $\mathbf{v}=(1,1,1)$.
  \item Compute $f_x$ and $f_y$ for $f(x,y)=(\ln x)(\ln y)$ with
        $x,y>1$.
  \item For $F(K,L)=K+2K^{1/2}L^{1/2}+L$ (which equals
        $(K^{1/2}+L^{1/2})^2$), compute $F_K, F_L$ directly and confirm
        agreement with the CES form.
  \item Compute the directional derivative of $f(x,y)=x^2y+y^3$ at
        $(-1,2)$ in the direction $\mathbf{v}=(4,-3)$.
  \item For $U(x,y)=\alpha\ln x+\beta\ln y$ (log-Cobb--Douglas), derive
        $\mathrm{MRS}_{xy}$ and evaluate at $(x,y)=(\alpha,\beta)$.
  \item Compute $\nabla f$ for $f(x,y)=(x^2+y^2)^{-1/2}$ (Newtonian
        potential in 2D).
  \item Compute $f_x$ and $f_y$ for $f(x,y)=\arctan(xy)$.
  \item For $F(K,L)=(K^\rho+L^\rho)^{1/\rho}$ with $\rho<1$, $\rho\neq0$
        (CES with constant returns), compute $F_K$ and show
        $F_K=(K/F)^{\rho-1}$.
  \item Compute the directional derivative of $f(x,y,z)=xyz$ at
        $(2,3,-1)$ in the direction $\mathbf{v}=(2,-1,2)$.
\end{problems}

\section{Solutions}

\begin{solutions}
  \item $f_x=20x^3-4xy^3$, $f_y=-6x^2y^2+5y^4$.

  \item $F_K=12\cdot\frac14K^{-3/4}L^{3/4}=3(L/K)^{3/4}$;
        $F_L=12\cdot\frac34K^{1/4}L^{-1/4}=9(K/L)^{1/4}$. At $(16,81)$:
        $(L/K)^{3/4}=(81/16)^{3/4}=27/8$ so $F_K=81/8$;
        $(K/L)^{1/4}=(16/81)^{1/4}=2/3$ so $F_L=6$.

  \item $f_x=2xe^{3y}$, $f_y=3x^2e^{3y}$.

  \item $\nabla f=(12x^2y-y^4,\ 4x^3-4xy^3)^T$.

  \item $f_x=2x-4y$, $f_y=-4x+3y^2$. At $(2,1)$: $\nabla f=(0,-5)^T$.
        Normalise: $\|\mathbf{v}\|=\sqrt{25+144}=13$,
        $\mathbf{u}=(5/13,-12/13)^T$. Then
        \[
          D_{\mathbf{u}}f=0\cdot\frac{5}{13}+(-5)\cdot
            \left(-\frac{12}{13}\right)=\frac{60}{13} .
        \]

  \item $F_K=6\cdot\frac23K^{-1/3}L^{1/3}=4(L/K)^{1/3}$,
        $F_L=2(K/L)^{2/3}$. At $(27,8)$: $(L/K)^{1/3}=2/3$ so
        $\mathrm{MP}_K=8/3$; $(K/L)^{2/3}=(27/8)^{2/3}=9/4$ so
        $\mathrm{MP}_L=9/2$.

  \item Chain rule with $u=3x-2y$: $f_x=4(3x-2y)^3\cdot3=12(3x-2y)^3$,
        $f_y=-8(3x-2y)^3$.

  \item $f_x=y^2e^{xy^2}$, $f_y=2xye^{xy^2}$.

  \item $F_K=15K^{-0.4}L^{0.5}$, $F_L=12.5K^{0.6}L^{-0.5}$.

  \item $f_x=2x-z^2e^x$, $f_y=3z$, $f_z=3y-2ze^x$. So
        $\nabla f=(2x-z^2e^x,\ 3z,\ 3y-2ze^x)^T$.

  \item $f_x=e^x\ln(1+y)$, $f_y=e^x/(1+y)$. At $(0,0)$:
        $\nabla f=(0,1)^T$. Normalise: $\|\mathbf{v}\|=\sqrt5$,
        $\mathbf{u}=(1/\sqrt5,2/\sqrt5)^T$. So
        \[
          D_{\mathbf{u}}f=0\cdot\frac{1}{\sqrt5}+1\cdot\frac{2}{\sqrt5}
            =\frac{2\sqrt5}{5} .
        \]

  \item $U_x=2/x$, $U_y=5/y$. $\mathrm{MRS}_{xy}=U_x/U_y
        =(2/x)/(5/y)=2y/(5x)$. At $(3,5)$: $\mathrm{MRS}=10/15=2/3$.

  \item For $x+y\neq0$: $f(x,y)=x-y$, so $f_x=1$, $f_y=-1$.

  \item With $u=x^2y+1$: $f_x=2xy/(x^2y+1)$, $f_y=x^2/(x^2y+1)$.

  \item $U_{x_1}=\frac12x_1^{-1/2}x_2^{1/3}x_3^{1/6}$,
        $U_{x_2}=\frac13x_1^{1/2}x_2^{-2/3}x_3^{1/6}$,
        $U_{x_3}=\frac16x_1^{1/2}x_2^{1/3}x_3^{-5/6}$. Equivalently
        $U_{x_i}=(\alpha_i/x_i)U$ with $\alpha=(1/2,1/3,1/6)$.

  \item $f_x=y$, $f_y=x-2z$, $f_z=-2y+2z$. At $(1,2,3)$:
        $\nabla f=(2,-5,2)^T$. Normalise: $\|\mathbf{v}\|=\sqrt{1+4+4}=3$,
        $\mathbf{u}=(1/3,2/3,2/3)^T$. Then
        $D_{\mathbf{u}}f=(2-10+4)/3=-4/3$.

  \item $F_K=0.25K^{-0.75}L^{0.75}$, $F_L=0.75K^{0.25}L^{-0.25}$. So
        $\mathrm{MRTS}_{LK}=F_L/F_K=3K/L$. At $(16,4)$:
        $\mathrm{MRTS}=48/4=12$.

  \item Expand: $f=x^3-3x^2y^2+xy-3y^3$. Hence $f_x=3x^2-6xy^2+y$,
        $f_y=-6x^2y+x-9y^2$.

  \item With $u=-y/x$: $\partial u/\partial x=y/x^2$,
        $\partial u/\partial y=-1/x$. Product rule:
        \[
          f_x=e^{-y/x}+x\cdot e^{-y/x}\cdot\frac{y}{x^2}
            =e^{-y/x}\left(1+\frac{y}{x}\right), \qquad
          f_y=x\cdot e^{-y/x}\cdot\left(-\frac1x\right)=-e^{-y/x} .
        \]

  \item Let $g=3K^{1/2}+2L^{1/2}$, so $F=g^2$. Then
        \[
          F_K=2g\cdot\frac32K^{-1/2}=\frac{3g}{K^{1/2}}
            =\frac{3(3K^{1/2}+2L^{1/2})}{K^{1/2}}, \qquad
          F_L=2g\cdot L^{-1/2}=\frac{2(3K^{1/2}+2L^{1/2})}{L^{1/2}} .
        \]

  \item With $u=y/x$, $f=\arctan u$, $f'=1/(1+u^2)$:
        \[
          f_x=\frac{1}{1+y^2/x^2}\cdot\left(-\frac{y}{x^2}\right)
            =\frac{-y}{x^2+y^2}, \qquad
          f_y=\frac{x}{x^2+y^2} .
        \]
        So $\nabla f=(-y/(x^2+y^2),\ x/(x^2+y^2))^T$.

  \item $f_K=\frac12K^{-1/2}L^{1/2}$, $f_L=\frac12K^{1/2}L^{-1/2}$. At
        $(4,9)$: $\nabla f=(\frac12\cdot\frac32,\ \frac12\cdot\frac23)^T
        =(3/4,1/3)^T$. Then $\|\nabla f\|^2=9/16+1/9=97/144$, so the
        maximum rate is $\sqrt{97}/12\approx0.821$ in direction
        $\mathbf{u}^*=(9,4)^T/\sqrt{97}$.

  \item $F=(K^{1/2}+L^{1/2})^2$. Replacing $(K,L)\to(\lambda K,\lambda L)$
        gives $\lambda F$, so $r=1$. With $g=K^{1/2}+L^{1/2}$:
        $F_K=2g\cdot\frac12K^{-1/2}=g/K^{1/2}$, $F_L=g/L^{1/2}$. Hence
        \[
          KF_K+LF_L=K^{1/2}g+L^{1/2}g=g(K^{1/2}+L^{1/2})=g^2=F. \checkmark
        \]

  \item Quotient rule, $D=x^2+y^2$:
        \[
          f_x=\frac{2D-(2x+3y)(2x)}{D^2}=\frac{-2x^2-6xy+2y^2}{(x^2+y^2)^2},
          \qquad f_y=\frac{3x^2-4xy-3y^2}{(x^2+y^2)^2} .
        \]

  \item Treat denominator as constant for $f_x$: $f_x=e^x/(1+e^y)$. For
        $f_y$:
        \[
          f_y=e^x\cdot\frac{-e^y}{(1+e^y)^2}=-\frac{e^{x+y}}{(1+e^y)^2} .
        \]

  \item Let $S=\alpha K^\rho+(1-\alpha)L^\rho$, $F=AS^{\nu/\rho}$.
        Treating $K$ constant, $\partial S/\partial L=\rho(1-\alpha)L^{\rho-1}$:
        \[
          F_L=A\cdot\frac{\nu}{\rho}S^{\nu/\rho-1}\cdot\rho(1-\alpha)L^{\rho-1}
            =\nu(1-\alpha)AL^{\rho-1}S^{\nu/\rho-1} .
        \]

  \item $f=\ln(x^2+y^2)$, $f_x=2x/(x^2+y^2)$, $f_y=2y/(x^2+y^2)$. At
        $(3,4)$: $x^2+y^2=25$, $\nabla f=(6/25,8/25)^T$. Normalise:
        $\|\mathbf{v}\|=\sqrt2$, $\mathbf{u}=(-1/\sqrt2,1/\sqrt2)^T$. Then
        \[
          D_{\mathbf{u}}f=\frac{6}{25}\cdot\left(-\frac{1}{\sqrt2}\right)
            +\frac{8}{25}\cdot\frac{1}{\sqrt2}=\frac{2}{25\sqrt2}
            =\frac{\sqrt2}{25} .
        \]

  \item $f_x=2xyz^3-4y^2z$, $f_y=x^2z^3-8xyz$,
        $f_z=3x^2yz^2-4xy^2+5$.

  \item Chain rule, $u=4x^2+9y^2$:
        \[
          f_x=\frac{8x}{2\sqrt u}=\frac{4x}{\sqrt{4x^2+9y^2}}, \qquad
          f_y=\frac{9y}{\sqrt{4x^2+9y^2}} .
        \]

  \item With $g=K^{1/2}+L^{1/2}$, $F=g^4$,
        $F_K=4g^3\cdot\frac12K^{-1/2}=2g^3/K^{1/2}$. At $(4,9)$: $g=5$,
        $K^{1/2}=2$, so $F_K=2\cdot125/2=125$.

  \item Let $r^2=x^2+y^2+z^2$. Then $f=\ln(r^2)$ and $f_{x_i}=2x_i/r^2$.
        So $\nabla f=(2/(x^2+y^2+z^2))(x,y,z)^T$.

  \item Use $\ln\sqrt{u}=\frac12\ln u$: $f=\frac12\ln(x^2+y^2)$. So
        $f_x=x/(x^2+y^2)$, $f_y=y/(x^2+y^2)$.

  \item $\nabla C=(4,9)^T$. With $Q=x^{1/2}y^{1/2}$: $Q_x=\frac12(y/x)^{1/2}$,
        $Q_y=\frac12(x/y)^{1/2}$, so
        $\nabla Q=\frac12\left((y/x)^{1/2},(x/y)^{1/2}\right)^T$. The
        condition $\nabla C=\lambda\nabla Q$ requires
        $4/9=Q_x/Q_y=y/x$, i.e.\ $y=4x/9$. Along this locus
        $\nabla Q\propto(2/3,3/2)^T\propto(4,9)^T=\nabla C$. \checkmark

  \item Chain rule: $f_x=5(1+xy)^4\cdot y=5y(1+xy)^4$,
        $f_y=5x(1+xy)^4$.

  \item $f_x=3x^2-2y$, $f_y=-2x+2y$. At $(1,2)$:
        $\nabla f=(3-4,-2+4)^T=(-1,2)^T$. Normalise: $\|\mathbf{v}\|=5$,
        $\mathbf{u}=(3/5,4/5)^T$. Then
        \[
          D_{\mathbf{u}}f=-1\cdot\frac35+2\cdot\frac45
            =\frac{-3+8}{5}=1 .
        \]

  \item $F_K=10\cdot0.3K^{-0.7}L^{0.7}=3(L/K)^{0.7}$;
        $F_L=7(K/L)^{0.3}$. At $K=L=100$: $F_K=3$, $F_L=7$. And
        $F=10\cdot100^{0.3}\cdot100^{0.7}=10\cdot100=1000$. Then
        $KF_K+LF_L=100\cdot3+100\cdot7=1000=F$. \checkmark (homogeneity
        degree $r=0.3+0.7=1$).

  \item $f_x=y\cos(xy)-\sin(x-y)$, $f_y=x\cos(xy)+\sin(x-y)$.

  \item $f_x=2xy-y^2+3$, $f_y=x^2-2xy-2$. At $(1,-1)$: $f_x=-2-1+3=0$,
        $f_y=1+2-2=1$. So $\nabla f(1,-1)=(0,1)^T$.

  \item $U_x=\alpha x^{\alpha-1}y^{1-\alpha}$,
        $U_y=(1-\alpha)x^\alpha y^{-\alpha}$. So
        \[
          \mathrm{MRS}_{xy}=\frac{U_x}{U_y}
            =\frac{\alpha x^{\alpha-1}y^{1-\alpha}}{(1-\alpha)x^\alpha y^{-\alpha}}
            =\frac{\alpha}{1-\alpha}\cdot\frac{y}{x} . \checkmark
        \]

  \item $f_x=ye^{xy}$, $f_y=xe^{xy}$. At $(1,0)$: $\nabla f=(0,1)^T$.
        Normalise: $\|\mathbf{v}\|=5$, $\mathbf{u}=(-3/5,4/5)^T$. Then
        $D_{\mathbf{u}}f=0\cdot(-3/5)+1\cdot(4/5)=4/5$.

  \item Write $f=e^{y\ln x}$. Then $f_x=e^{y\ln x}\cdot y/x=yx^{y-1}$ and
        $f_y=e^{y\ln x}\cdot\ln x=x^y\ln x$.

  \item On $\{2K<3L\}$ the binding factor is capital: $F=2K$, so
        $F_K=2$, $F_L=0$. On $\{2K>3L\}$ labour binds: $F=3L$, so
        $F_K=0$, $F_L=3$. (Non-differentiable on $\{2K=3L\}$.)

  \item $f_x=yz+2x$, $f_y=xz$, $f_z=xy-2z$. At $(2,-1,3)$:
        $f_x=-3+4=1$, $f_y=6$, $f_z=-2-6=-8$. So
        $\nabla f=(1,6,-8)^T$.

  \item $\nabla f=(2x,8y)^T$. At $(3,1)$: $\nabla f=(6,8)^T$. Steepest
        descent is $-\nabla f/\|\nabla f\|$. $\|\nabla f\|=10$, so unit
        descent direction is $(-3/5,-4/5)^T$, with rate of decrease
        $-\|\nabla f\|=-10$ (i.e.\ $f$ decreases at rate 10).

  \item $f_x=e^y-ye^x$, $f_y=xe^y-e^x$.

  \item Quotient rule with $D=x^2+y^2$:
        \[
          f_x=\frac{D-x\cdot2x}{D^2}=\frac{y^2-x^2}{(x^2+y^2)^2}, \qquad
          f_y=\frac{-x\cdot2y}{D^2}=\frac{-2xy}{(x^2+y^2)^2} .
        \]

  \item $U_x=1$, $U_y=1/\sqrt{y}$. So $\mathrm{MRS}_{xy}=U_x/U_y
        =\sqrt{y}$ (depends only on $y$, quasi-linear utility).

  \item $f_x=2x$, $f_y=2y$, $f_z=-2z$. At $(2,-1,2)$:
        $\nabla f=(4,-2,-4)^T$. Normalise: $\|\mathbf{v}\|=\sqrt{4+4+1}=3$,
        $\mathbf{u}=(2/3,2/3,1/3)^T$. Then
        \[
          D_{\mathbf{u}}f=4\cdot\frac23+(-2)\cdot\frac23+(-4)\cdot\frac13
            =\frac{8-4-4}{3}=0 .
        \]
        (We are moving tangent to the level surface.)

  \item $f_x=\sec^2(x+y)+\sec^2(x-y)$, $f_y=\sec^2(x+y)-\sec^2(x-y)$.

  \item Let $S=\alpha K+\beta L$, $F=S^r$. Then $F_K=rS^{r-1}\alpha$,
        $F_L=rS^{r-1}\beta$. Hence
        \[
          KF_K+LF_L=rS^{r-1}(\alpha K+\beta L)=rS^{r-1}\cdot S=rS^r=rF.
            \checkmark
        \]

  \item $\nabla f=(-e^{-x}\sin y+e^x\cos y,\ e^{-x}\cos y-e^x\sin y)^T$.

  \item $f_x=x/\sqrt{x^2+y^2}$, $f_y=y/\sqrt{x^2+y^2}$. At $(3,4)$:
        $\sqrt{x^2+y^2}=5$, $\nabla f=(3/5,4/5)^T$. The direction
        $\mathbf{v}=(1,0)$ is already unit. So
        $D_{\mathbf{u}}f=3/5\cdot1+4/5\cdot0=3/5$.

  \item On $\{2x<3y\}$: $U=2x$, so $U_x=2$, $U_y=0$. On $\{2x>3y\}$:
        $U=3y$, so $U_x=0$, $U_y=3$. (Non-differentiable on the ray
        $2x=3y$.)

  \item Let $L=\ln(x+y)$, $D=x-y$. By the quotient rule:
        \[
          f_x=\frac{(1/(x+y))(x-y)-\ln(x+y)\cdot1}{(x-y)^2}
            =\frac{1/(x+y)\cdot(x-y)-\ln(x+y)}{(x-y)^2} .
        \]
        Simplifying
        $f_x=\dfrac{(x-y)/(x+y)-\ln(x+y)}{(x-y)^2}$. Similarly,
        \[
          f_y=\frac{(1/(x+y))(x-y)-\ln(x+y)\cdot(-1)}{(x-y)^2}
            =\frac{(x-y)/(x+y)+\ln(x+y)}{(x-y)^2} .
        \]

  \item $F_K=A\alpha K^{\alpha-1}L^\beta$, $F_L=A\beta K^\alpha L^{\beta-1}$.
        So
        \[
          \mathrm{MRTS}_{LK}=\frac{F_L}{F_K}=\frac{\beta}{\alpha}\cdot
            \frac{K}{L} .
        \]
        This depends only on the ratio $K/L$ and the exponent ratio
        $\beta/\alpha$: it is scale-invariant.

  \item $\nabla f=(2x,-2y)^T$. At $(2,-1)$: $\nabla f=(4,2)^T$. Direction
        of steepest ascent: $\mathbf{u}^*=(4,2)/\sqrt{20}=(2,1)/\sqrt5$.
        Maximum rate $=\|\nabla f\|=\sqrt{20}=2\sqrt5$.

  \item $f_x=\sinh(x+y)+\cosh(x-y)$, $f_y=\sinh(x+y)-\cosh(x-y)$.

  \item $\nabla f=(e^{x+y+z},e^{x+y+z},e^{x+y+z})^T$. At $(0,0,0)$:
        $\nabla f=(1,1,1)^T$. Normalise $\mathbf{v}=(1,1,1)$:
        $\|\mathbf{v}\|=\sqrt3$, $\mathbf{u}=(1,1,1)/\sqrt3$. Then
        $D_{\mathbf{u}}f=(1+1+1)/\sqrt3=\sqrt3$.

  \item $f_x=(\ln y)/x$, $f_y=(\ln x)/y$.

  \item $F_K=1+L^{1/2}/K^{1/2}$, $F_L=1+K^{1/2}/L^{1/2}$. From
        $F=(K^{1/2}+L^{1/2})^2$ with $g=K^{1/2}+L^{1/2}$:
        $F_K=g/K^{1/2}=(K^{1/2}+L^{1/2})/K^{1/2}=1+L^{1/2}/K^{1/2}$.
        \checkmark Same for $F_L$.

  \item $f_x=2xy$, $f_y=x^2+3y^2$. At $(-1,2)$: $\nabla f=(-4,1+12)^T
        =(-4,13)^T$. Normalise: $\|\mathbf{v}\|=5$,
        $\mathbf{u}=(4/5,-3/5)^T$. Then
        \[
          D_{\mathbf{u}}f=-4\cdot\frac45+13\cdot\left(-\frac35\right)
            =-\frac{16+39}{5}=-11 .
        \]

  \item $U_x=\alpha/x$, $U_y=\beta/y$.
        $\mathrm{MRS}_{xy}=\dfrac{\alpha/x}{\beta/y}=\dfrac{\alpha y}{\beta x}$.
        At $(\alpha,\beta)$: $\mathrm{MRS}=(\alpha\cdot\beta)/(\beta\cdot\alpha)=1$.

  \item $f=(x^2+y^2)^{-1/2}$. Chain rule:
        \[
          f_x=-\frac12(x^2+y^2)^{-3/2}\cdot2x=\frac{-x}{(x^2+y^2)^{3/2}},
          \qquad f_y=\frac{-y}{(x^2+y^2)^{3/2}} .
        \]
        So $\nabla f=-(x,y)^T/(x^2+y^2)^{3/2}$.

  \item Chain rule with $u=xy$:
        \[
          f_x=\frac{y}{1+x^2y^2}, \qquad f_y=\frac{x}{1+x^2y^2} .
        \]

  \item Let $S=K^\rho+L^\rho$, $F=S^{1/\rho}$. Then
        $F_K=\frac1\rho S^{1/\rho-1}\cdot\rho K^{\rho-1}
        =K^{\rho-1}S^{1/\rho-1}$. Since $S^{1/\rho}=F$, we have $S=F^\rho$
        and so $S^{1/\rho-1}=F^{1-\rho}\cdot F^{-1}\cdot F=F^{1-\rho}$.
        Thus $F_K=K^{\rho-1}F^{1-\rho}=(F/K)^{1-\rho}=(K/F)^{\rho-1}$.
        \checkmark

  \item $f_x=yz$, $f_y=xz$, $f_z=xy$. At $(2,3,-1)$:
        $\nabla f=(3\cdot(-1),2\cdot(-1),2\cdot3)^T=(-3,-2,6)^T$.
        Normalise $\mathbf{v}=(2,-1,2)$: $\|\mathbf{v}\|=3$,
        $\mathbf{u}=(2/3,-1/3,2/3)^T$. Then
        \[
          D_{\mathbf{u}}f=-3\cdot\frac23+(-2)\cdot\left(-\frac13\right)
            +6\cdot\frac23=-2+\frac23+4=2+\frac23=\frac83 .
        \]
\end{solutions}

\section{Additional Problems}

\begin{problems}
  \item Compute $f_x$ and $f_y$ for $f(x,y)=3x^5-2x^3y^2+4y^4$.
  \item For $F(K,L)=20\,K^{1/3}L^{2/3}$, find $F_K$ and $F_L$ at
        $(K,L)=(8,27)$.
  \item Compute $f_x$ and $f_y$ for $f(x,y)=e^{xy}$.
  \item Compute $\nabla f$ for $f(x,y)=2x^2y-3xy^3+5$.
  \item Compute the directional derivative of $f(x,y)=x^2+2xy-y^2$ at
        $(1,2)$ in the direction $\mathbf{v}=(3,4)$.
  \item For $F(K,L)=8K^{3/4}L^{1/4}$, compute $\mathrm{MP}_K$ and
        $\mathrm{MP}_L$ at $(K,L)=(16,16)$.
  \item Compute $f_x$ and $f_y$ for $f(x,y)=\ln(2x+3y)$.
  \item For $U(x,y)=3\ln(x)+4\ln(y)$, derive $\mathrm{MRS}_{xy}$ and
        evaluate at $(2,3)$.
  \item Compute $f_x$ and $f_y$ for $f(x,y)=(x+y^2)^4$.
  \item Compute $\nabla f$ for $f(x,y,z)=xy+yz+zx$.
  \item For a CES production $F(K,L)=5(K^{1/3}+L^{1/3})^3$, find $F_K$.
  \item Compute the directional derivative of $f(x,y)=x^2y$ at $(2,1)$
        in the direction $\mathbf{v}=(1,-1)$.
  \item For $U(x,y)=x^\alpha y^\beta$ (general Cobb--Douglas,
        $\alpha,\beta>0$), derive $\mathrm{MRS}_{xy}$ and show it depends
        only on the ratio $y/x$.
  \item Compute $f_x$ and $f_y$ for $f(x,y)=\dfrac{x-y}{x+y}$.
  \item For $f(K,L)=K^{1/2}L^{1/2}$ at $(K,L)=(16,9)$, find the unit
        direction of steepest ascent and the maximum rate of increase.
  \item Classify the open positive orthant
        $Q=\{(x,y)\in\R^2 : x>0,\ y>0\}$: is it open, closed, or
        neither? Is it bounded? Is it compact?
  \item Classify the budget set
        $B=\{(x,y)\in\R^2 : 2x+3y\le12,\ x\ge0,\ y\ge0\}$: open,
        closed, or neither? Bounded? Compact?
  \item Classify $M=\{(x,y)\in\R^2 : xy\ge1,\ x>0\}$: open, closed, or
        neither? Bounded? Compact? \emph{Hint:} check whether a
        sequence in $M$ can converge to a point with $x=0$.
  \item Classify the open disc
        $D=\{(x,y)\in\R^2 : x^2+y^2<4\}$: open, closed, or neither?
        Bounded? Compact?
  \item Classify (a) the half-open strip
        $S=\{(x,y)\in\R^2 : 0\le x<1\}$ and (b) the finite set
        $E=\{(0,0),\ (1,2),\ (3,1)\}$: open, closed, or neither?
        Bounded? Compact?
  \item Does the Weierstrass theorem guarantee a maximiser for the
        continuous function $u(x,y)=x^{1/2}y^{1/2}$ (a) on the budget
        set $B=\{(x,y): 2x+3y\le12,\ x\ge0,\ y\ge0\}$, (b) on the open
        positive orthant $Q=\{(x,y): x>0,\ y>0\}$? Justify each
        answer.
\end{problems}

\section{Additional Solutions}

\begin{solutions}
  \item $f_x=15x^4-6x^2y^2$, $f_y=-4x^3y+16y^3$.

  \item $F_K=20\cdot\frac13K^{-2/3}L^{2/3}=\frac{20}{3}(L/K)^{2/3}$;
        $F_L=20\cdot\frac23K^{1/3}L^{-1/3}=\frac{40}{3}(K/L)^{1/3}$. At
        $(8,27)$: $(L/K)^{2/3}=(27/8)^{2/3}=9/4$, so $F_K=15$;
        $(K/L)^{1/3}=(8/27)^{1/3}=2/3$, so
        $F_L=\frac{40}{3}\cdot\frac23=\frac{80}{9}$.

  \item $f_x=ye^{xy}$, $f_y=xe^{xy}$.

  \item $\nabla f=(4xy-3y^3,\ 2x^2-9xy^2)^T$.

  \item $f_x=2x+2y$, $f_y=2x-2y$. At $(1,2)$: $\nabla f=(6,-2)^T$.
        Normalise: $\|\mathbf{v}\|=5$, $\mathbf{u}=(3/5,4/5)^T$. Then
        \[
          D_{\mathbf{u}}f=6\cdot\frac35+(-2)\cdot\frac45
            =\frac{18-8}{5}=2 .
        \]

  \item $F_K=8\cdot\frac34K^{-1/4}L^{1/4}=6(L/K)^{1/4}$,
        $F_L=2(K/L)^{3/4}$. At $(16,16)$: $\mathrm{MP}_K=6$,
        $\mathrm{MP}_L=2$.

  \item $f_x=\dfrac{2}{2x+3y}$, $f_y=\dfrac{3}{2x+3y}$.

  \item $U_x=3/x$, $U_y=4/y$.
        $\mathrm{MRS}_{xy}=U_x/U_y=\dfrac{3y}{4x}$. At $(2,3)$:
        $\mathrm{MRS}=\dfrac98$.

  \item Chain: $f_x=4(x+y^2)^3$,
        $f_y=4(x+y^2)^3\cdot2y=8y(x+y^2)^3$.

  \item $f_x=y+z$, $f_y=x+z$, $f_z=x+y$. So
        $\nabla f=(y+z,\ x+z,\ x+y)^T$.

  \item Let $g=K^{1/3}+L^{1/3}$, so $F=5g^3$. Then
        $F_K=5\cdot3g^2\cdot\frac13K^{-2/3}=5g^2K^{-2/3}
        =\dfrac{5(K^{1/3}+L^{1/3})^2}{K^{2/3}}$.

  \item $f_x=2xy$, $f_y=x^2$. At $(2,1)$: $\nabla f=(4,4)^T$. Normalise:
        $\|\mathbf{v}\|=\sqrt2$, $\mathbf{u}=(1/\sqrt2,-1/\sqrt2)^T$.
        Then
        \[
          D_{\mathbf{u}}f=4\cdot\frac{1}{\sqrt2}
            +4\cdot\left(-\frac{1}{\sqrt2}\right)=0 .
        \]
        (Tangent to the level curve.)

  \item $U_x=\alpha x^{\alpha-1}y^\beta$,
        $U_y=\beta x^\alpha y^{\beta-1}$. So
        \[
          \mathrm{MRS}_{xy}=\frac{U_x}{U_y}
            =\frac{\alpha}{\beta}\cdot\frac{y}{x} .
        \]
        Depends only on $y/x$ (and the exponent ratio), not on the
        scale. \checkmark

  \item Quotient rule with $D=x+y$:
        \[
          f_x=\frac{D-(x-y)}{D^2}=\frac{2y}{(x+y)^2}, \qquad
          f_y=\frac{-D-(x-y)}{D^2}=\frac{-2x}{(x+y)^2} .
        \]

  \item $f_K=\frac12K^{-1/2}L^{1/2}$, $f_L=\frac12K^{1/2}L^{-1/2}$. At
        $(16,9)$: $f_K=\frac12\cdot\frac34=3/8$;
        $f_L=\frac12\cdot\frac43=2/3$. So $\nabla f=(3/8,2/3)^T$.
        Magnitude:
        $\|\nabla f\|=\sqrt{9/64+4/9}=\sqrt{81/576+256/576}
        =\sqrt{337/576}=\sqrt{337}/24\approx0.765$. Unit direction of
        steepest ascent:
        $\mathbf{u}^*=(3/8,2/3)^T/(\sqrt{337}/24)=(9,16)^T/\sqrt{337}$.

  \item \textbf{Open.} For any $(x,y)\in Q$ the ball of radius
        $r=\min\{x,y\}>0$ around $(x,y)$ stays inside $Q$, so every
        point is interior. \textbf{Not closed:} the boundary points
        $(x,0)$ and $(0,y)$ with $x,y\ge0$ (the coordinate half-axes,
        including the origin) do not belong to $Q$. \textbf{Not
        bounded:} $Q$ contains $(t,t)$ for every $t>0$, and
        $\|(t,t)\|=t\sqrt2\to\infty$. Hence $Q$ is open, unbounded and
        \textbf{not compact}.

  \item \textbf{Closed.} $B$ is the intersection of the three sets
        $\{2x+3y\le12\}$, $\{x\ge0\}$ and $\{y\ge0\}$, each defined by
        a non-strict inequality on a continuous function; each is
        closed, and an intersection of closed sets is closed (every
        boundary point of $B$ satisfies all three inequalities, so it
        lies in $B$). \textbf{Not open:} $(0,0)\in B$, but every ball
        around the origin contains points with negative coordinates,
        which are outside $B$; so $(0,0)$ is not interior.
        \textbf{Bounded:} from the constraints, $0\le x\le6$ and
        $0\le y\le4$, so $B\subseteq B_8(\mathbf{0})$. Closed and
        bounded in $\R^2$: $B$ is \textbf{compact}.

  \item First note the constraint $x>0$ together with $xy\ge1$ forces
        $y\ge1/x>0$. \textbf{Closed:} suppose
        $(x_n,y_n)\in M$ with $(x_n,y_n)\to(a,b)$. If $a=0$, then
        $y_n\to b$ finite would give $x_ny_n\to0<1$, contradicting
        $x_ny_n\ge1$; so $a>0$, and passing to the limit in
        $x_ny_n\ge1$ gives $ab\ge1$, hence $(a,b)\in M$. $M$ contains
        all its limit (boundary) points, so it is closed: even
        though it is written with a strict inequality, that inequality
        is never active. \textbf{Not open:} $(1,1)\in M$ lies on the
        hyperbola $xy=1$, and every ball around it contains points
        with $xy<1$. \textbf{Not bounded:} $(t,t)\in M$ for all
        $t\ge1$. Hence $M$ is closed, unbounded and \textbf{not
        compact}.

  \item \textbf{Open.} For $(x,y)\in D$ put
        $r=2-\sqrt{x^2+y^2}>0$; the ball $B_r(x,y)$ lies in $D$ by the
        triangle inequality. \textbf{Not closed:} the boundary circle
        $\{x^2+y^2=4\}$ is not contained in $D$ (e.g.\ $(2,0)\notin
        D$). \textbf{Bounded:} $D\subseteq B_3(\mathbf{0})$. Bounded
        but not closed: $D$ is \textbf{not compact}.

  \item (a) \textbf{Neither open nor closed.} The points $(0,y)$
        belong to $S$ but are boundary points (every ball around them
        contains points with $x<0$), so $S$ is not open; the boundary
        points $(1,y)$ do not belong to $S$, so $S$ is not closed.
        \textbf{Unbounded:} $(0,t)\in S$ for all $t$. Hence
        \textbf{not compact}. (b) \textbf{Closed:} around any point
        not in $E$ there is a ball avoiding all three points (a
        finite set always contains all its boundary points:
        indeed $E$ is its own boundary). \textbf{Not open:} no ball
        around $(0,0)$ lies inside $E$. \textbf{Bounded:}
        $E\subseteq B_4(\mathbf{0})$. Closed and bounded: $E$ is
        \textbf{compact}.

  \item (a) \textbf{Yes.} $B$ is compact (closed and bounded, see the
        budget-set problem above) and $u$ is continuous on $B$, so by
        the Weierstrass theorem $u$ attains both a maximum and a
        minimum on $B$; a maximiser exists. (b) \textbf{No.} $Q$ is
        not compact (it is neither closed nor bounded), so Weierstrass
        does not apply: and indeed no maximiser exists:
        $u(t,t)=t\to\infty$ as $t\to\infty$, so $u$ is unbounded above
        on $Q$.
\end{solutions}
